\documentclass{standalone}
\usepackage{standalone}

\begin{document}
\chapter{System Design and Architecture}

\section{High-Level Architecture}
Sarabo has two parts: a client that runs in the browser and an API server. They rely on outside services for identity, payment, storage, email and AI (Figure~\ref{fig:arch}). Both parts run on Vercel. The API is a serverless function in the Mumbai region, close to the MongoDB Atlas cluster that holds all the data, and Vercel also triggers a job once a day.

\begin{figure}[htbp]
\centering
\begin{tikzpicture}[font=\small]
  \node[hl,minimum width=3.2cm] (client) {React 19 + Vite\\Web client\\{\scriptsize Vercel (static)}};
  \node[hl,minimum width=3.2cm,right=2.4cm of client] (api) {Express 5 REST API\\{\scriptsize Vercel serverless, Mumbai}};
  \node[db,right=1.6cm of api] (db) {MongoDB\\Atlas};
  \node[box,below=1.9cm of client,xshift=-1.2cm] (fb) {Firebase\\Authentication};
  \node[box,right=4mm of fb] (stripe) {Stripe\\Checkout, webhook,\\refunds};
  \node[box,right=4mm of stripe] (supa) {Supabase Storage\\private bucket};
  \node[box,right=4mm of supa] (gem) {Google\\Gemini};
  \node[box,right=4mm of gem] (mail) {Gmail\\SMTP};
  \node[box,above=1.1cm of api] (cron) {Vercel Cron\\{\scriptsize daily 18:00 Dhaka}};
  \draw[arr,<->] (client) -- node[lbl,above]{HTTPS / JSON} node[lbl,below]{ID token} (api);
  \draw[arr,<->] (api) -- (db);
  \draw[arr] (client) -- node[lbl,left]{sign in} (fb);
  \draw[arr] (api) -- (fb.north east);
  \draw[arr,<->] (api) -- (stripe.north);
  \draw[arr,<->] (api) -- (supa.north);
  \draw[arr] (api) -- (gem.north);
  \draw[arr] (api) -- (mail.north);
  \draw[arr] (cron) -- (api);
  \draw[arr,dashed] (client.south east) -- node[lbl,pos=0.55]{signed PUT} (supa.north west);
\end{tikzpicture}
\caption{Sarabo high-level architecture}
\label{fig:arch}
\end{figure}

\section{Client Architecture}
The client has about 337 source files, organised into five areas:
\begin{enumerate}
  \item \textbf{Routing and guards}: the public pages (home, services, areas, jobs, tracking, about, become a technician), the sign-in pages, and a dashboard that \texttt{CustomerRoute}, \texttt{TechnicianRoute} and \texttt{AdminRoute} divide by role. The server checks every rule again on its side.
  \item \textbf{Feature components}: the request form and AI helper, photos, pickup slots, offers, inspection, quote, payment, repair, warranty, feedback, wallet, notifications, support chat and administration.
  \item \textbf{Server-state hooks}: TanStack Query with query keys kept in one place, which handles caching, loading and error states, mutations and refreshing only the data that changed.
  \item \textbf{Presentation utilities}: the ``service spine'' that maps each status to a stage, currency formatting, eligibility helpers and an error message for each action.
  \item \textbf{Design system}: shared design tokens, light and dark themes, and Radix UI primitives.
\end{enumerate}

\section{Server Architecture}
The server has about 191 JavaScript files, tests included. It is built in layers: \emph{routes} $\rightarrow$ \emph{middleware} $\rightarrow$ \emph{controllers} $\rightarrow$ \emph{services} / \emph{utilities} $\rightarrow$ \emph{database}. The middleware gives each request an ID, logs it safely, verifies the Firebase token and waits until the database is ready. Each controller is responsible for the state rules of one domain. Services hold the logic that several controllers share, such as eligibility, booking, refunds, settlement, notifications, matching, AI and email. External providers are wrapped in adapters, so the tests can swap them for fakes.

The Stripe webhook route is registered before the JSON parser, because checking the signature needs the raw request body. If the database goes down for a moment, the routes stay registered, so requests fail with a clear 503 error instead of a misleading 404.

\section{Data Collections}
The data is stored in 21 MongoDB collections, listed in Table~\ref{tab:collections}.

\begin{table}[htbp]
\centering
\caption{MongoDB collections}
\label{tab:collections}
\small
\begin{tabularx}{\textwidth}{@{}p{4.6cm}X@{}}
\toprule
\textbf{Collection} & \textbf{Purpose} \\
\midrule
\texttt{users} & Accounts and the authoritative role. \\
\texttt{technicians} & Application, approval status, expertise, region/district, work status, bKash details. \\
\texttt{repair\_requests} & The repair case with all lifecycle sub-documents. \\
\texttt{job\_applications} & Technician offers on open requests and their states. \\
\texttt{service\_definitions} & Versioned catalogue of services and prices. \\
\texttt{payments} & Stripe payment rows for inspection fees and repairs, including refund state. \\
\texttt{checkout\_sessions} & Checkout idempotency and expiry. \\
\texttt{technician\_withdrawals} & Withdrawal requests and admin processing. \\
\texttt{technician\_reviews} & Customer ratings with moderation history. \\
\texttt{technician\_reports} & Private customer reports with status and admin notes. \\
\texttt{notifications} & In-app notifications with de-duplication keys. \\
\texttt{tracking\_events} & Sanitized public timeline. \\
\texttt{damage\_upload\_sessions} & Customer photo upload sessions. \\
\texttt{repair\_evidence\_sessions} & Technician completion-evidence sessions. \\
\texttt{deletion\_cleanups} & Retry records for storage clean-up after deletion. \\
\texttt{turned\_away\_requests} & Hashed count of refused demand (180-day expiry). \\
\texttt{email\_log} & Sent/failed email record (90-day expiry). \\
\texttt{support\_rate\_limits} & Shared AI rate-limit counters (auto-expiring). \\
\texttt{ai\_report\_summaries} & Last AI summary per report period. \\
\texttt{warranty\_claims} & Warranty claims with their full decision and refund history. \\
\texttt{warranty\_photo\_uploads} & Short-lived upload sessions for warranty photos. \\
\bottomrule
\end{tabularx}
\end{table}

\section{Repair Request Document}
Each request has \texttt{schemaVersion: 2}, a \texttt{product} snapshot (category, brand and model) and its own copy of the chosen service with its price rule. Because of this copy, editing the catalogue later never changes a request that already exists (Figure~\ref{fig:requestdoc}).

\begin{figure}[htbp]
\centering
\begin{tikzpicture}[font=\small]
  \node[db,minimum width=2.8cm] (r) {repair\_requests};
  \node[box,above left=10mm and 18mm of r] (p) {product\\brand, model};
  \node[box,above=10mm of r] (s) {service +\\price snapshot};
  \node[box,above right=10mm and 18mm of r] (d) {damage\\images (max 3)};
  \node[box,left=12mm of r] (pk) {pickup slot};
  \node[box,right=12mm of r] (ip) {inspection fee\\payment / refund};
  \node[box,below left=10mm and 18mm of r] (i) {inspection};
  \node[box,below=10mm of r] (q) {quote +\\quoteHistory};
  \node[box,below right=10mm and 18mm of r] (rp) {payment, repair,\\evidence, settlement};
  \foreach \n in {p,s,d,pk,ip,i,q,rp} \draw[arr] (r) -- (\n);
\end{tikzpicture}
\caption{Conceptual structure of a V2 repair request}
\label{fig:requestdoc}
\end{figure}

\section{State Model}
The status field, \texttt{deliveryStatus}, and several of its values come from the parcel-delivery codebase that Sarabo started from. I kept them as a fixed internal contract, and users only ever see repair terms (Table~\ref{tab:states}).

\begin{table}[htbp]
\centering
\caption{Repair lifecycle states}
\label{tab:states}
\small
\begin{tabularx}{\textwidth}{@{}p{4.4cm}X@{}}
\toprule
\textbf{Stored state} & \textbf{Meaning shown to users} \\
\midrule
\texttt{pending-pickup} & Waiting for a technician; open on the jobs board. \\
\texttt{assignment\_pending} & Offered directly by an admin; technician to accept or reject. \\
\texttt{driver\_assigned} & Technician booked. \\
\texttt{rider\_arriving} & Technician on the way. \\
\texttt{parcel\_picked\_up} & Device collected. \\
\texttt{inspection\_completed} & Inspection submitted. \\
\texttt{quote\_submitted} & Final price waiting for the customer. \\
\texttt{quote\_approved} & Customer accepted the price. \\
\texttt{quote\_rejected} & Customer declined the price. \\
\texttt{payment\_completed} & Customer paid. \\
\texttt{repair\_in\_progress} & Repair underway. \\
\texttt{repair\_completed} & Repair finished; waiting for hand-over. \\
\texttt{parcel\_delivered} & Customer confirmed receipt (terminal). \\
\texttt{cancelled} & Ended through an allowed path. \\
\bottomrule
\end{tabularx}
\end{table}

Other records have their own, smaller sets of states. A job application (an offer) can be \texttt{pending}, \texttt{accepted}, \texttt{not\_chosen}, \texttt{withdrawn}, \texttt{no\_show} or \texttt{declined}. A settlement goes from \emph{pending} to \emph{available}, or to \emph{reversed} after a refund. A warranty claim starts as \emph{submitted}, goes through re-inspection and a decision, and ends as \emph{resolved} or \emph{refunded}. A withdrawal goes from \emph{requested} to either \emph{paid} or \emph{rejected}.

\section{The Service Spine}
That many states would be hard for a customer to follow, so the interface groups them into four stages that are easy to understand (Figure~\ref{fig:spine}). These stages appear on the home page, the tracking page and every dashboard. Confirming receipt is not a fifth stage. It is the hand-over action that comes after ``Repaired''.

\begin{figure}[htbp]
\centering
\begin{tikzpicture}[node distance=10mm]
  \node[hl,minimum width=2.4cm] (a) {Request};
  \node[hl,minimum width=2.4cm,right=of a] (b) {Inspect};
  \node[hl,minimum width=2.4cm,right=of b] (c) {Approve};
  \node[hl,minimum width=2.4cm,right=of c] (d) {Repaired};
  \draw[arr] (a)--(b); \draw[arr] (b)--(c); \draw[arr] (c)--(d);
\end{tikzpicture}
\caption{The four-stage service spine shown to users}
\label{fig:spine}
\end{figure}

\section{API Design}
The API has more than 120 endpoints, grouped into families by resource. Table~\ref{tab:api} shows some typical ones, and Appendix~\ref{app:api} gives a longer list.

\begin{table}[htbp]
\centering
\caption{Representative API endpoints}
\label{tab:api}
\footnotesize
\begin{tabularx}{\textwidth}{@{}p{1.2cm}p{8cm}X@{}}
\toprule
\textbf{Method} & \textbf{Route} & \textbf{Purpose} \\
\midrule
POST & \nolinkurl{/repair-requests} & Create a request \\
GET & \nolinkurl{/jobs} & Public open-jobs board \\
POST & \nolinkurl{/repair-requests/:id/applications} & Technician sends an offer \\
POST & \nolinkurl{/repair-requests/:id/applications/:aid/accept} & Customer chooses an offer \\
POST & \nolinkurl{/repair-requests/:id/inspection} & Submit inspection \\
POST & \nolinkurl{/repair-requests/:id/quote/decision} & Approve or reject quote \\
POST & \nolinkurl{/repair-requests/:id/checkout-session} & Start repair payment \\
POST & \nolinkurl{/stripe-webhook} & Signed payment events \\
POST & \nolinkurl{/repair-requests/:id/repair/complete} & Complete repair \\
POST & \nolinkurl{/repair-requests/:id/review} & Rate the technician \\
POST & \nolinkurl{/ai/request-draft} & AI request helper \\
GET & \nolinkurl{/admin/attention} & Items needing admin action \\
GET & \nolinkurl{/public/trackings/:code} & Sanitized tracking \\
\bottomrule
\end{tabularx}
\end{table}

\end{document}
